\subsection{Stochastic online tracking with an explicit perturbation budget}
\label{app:tracking}

The landscape result concerns a fixed covariance. During training both its
eigenvectors and eigenvalues change, and the tracked basis is driven by noisy
mini-batches. We first establish a finite-horizon bound on objective excess.
This avoids silently replacing PL by a two-trajectory contraction property.

\begin{assumption}[Bounded, separated tracking interval]
\label{ass:tracking}
On the interval considered, $C_t\succeq0$,
$\|C_t\|_{\mathrm{op}}\le C_0$, and
$\lambda_r(C_t)-\lambda_{r+1}(C_t)\ge\Delta_0>0$.
Every update remains in the region
$\theta_{\max}(U_t,U_t^\star)\le R_0<\pi/2$.
Its descent representation is
\begin{equation}
 U_{t+1}=\operatorname{Retr}_{U_t}\{-\rho_t(g_t+\zeta_t+b_t)\},
 \qquad g_t=\operatorname{grad}\Phi_t(U_t),\label{eq:tracking-update}
\end{equation}
where $\mathbb E_t\zeta_t=0$,
$\mathbb E_t\|\zeta_t\|^2\le\sigma_t^2$, and
$\mathbb E_t\|b_t\|^2\le B_t^2$.
The retraction pullback satisfies
\begin{equation}
 \Phi_t(\operatorname{Retr}_U\xi)
 \le\Phi_t(U)+\langle\operatorname{grad}\Phi_t(U),\xi\rangle
                    +\frac{L_R}{2}\|\xi\|^2\label{eq:pullback-smooth}
\end{equation}
for all actual steps, and
$\|C_{t+1}-C_t\|\le v_t$ almost surely, with deterministic $v_t$.
\end{assumption}

Proposition~\ref{prop:qr-calculus} verifies \eqref{eq:pullback-smooth}
with the sharp, dimension-free choice
$L_R=\sup_t[\lambda_1(C_t)-\lambda_d(C_t)]\le C_0$.
No small-angle assumption is required for this pullback inequality itself;
the principal-angle restriction is used only by the PL inequality.

The bias $b_t$ records deviations from a fresh stochastic tangent-gradient
step. It includes the Grassmann momentum lag, deterministic approximation
errors, and any clipping bias. It is permitted to correlate with $\zeta_t$.
This is useful for the implemented momentum, whose error cannot be treated
as independent zero-mean noise.

\begin{lemma}[Ky Fan completion of squares for moving objective excess]
\label{lem:excess-drift}
Let $C,C'$ be symmetric, let $C$ have rank-$r$ gap at least $\Delta_0>0$,
and put $v=\|C'-C\|_F$. No gap assumption on $C'$ is needed. For every
rank-$r$ projector $P$ and every $a>0$,
\begin{equation}
 h_{C'}(P)\le(1+a)h_C(P)+\frac{1+a^{-1}}{4\Delta_0}v^2.
 \label{eq:excess-drift}
\end{equation}
For fixed $a$, the coefficient $(1+a^{-1})/4$ is optimal to second order.
\end{lemma}
\begin{proof}
Let $P_\star,P_\star'$ maximize the respective Ky Fan traces
\cite{fan1951maximum}, $E=C'-C$, and $D=P_\star'-P_\star$. Then
\begin{align}
 h_{C'}(P)-h_C(P)
 &=\tfrac12\operatorname{tr}\{E(P_\star-P)\}
                      +\tfrac12\operatorname{tr}(C'D),\\
 \tfrac12\operatorname{tr}(C'D)
 &=-h_C(P_\star')+\tfrac12\langle E,D\rangle\\
 &\le-\tfrac14\Delta_0\|D\|_F^2+\tfrac12v\|D\|_F
 \le\frac{v^2}{4\Delta_0}.
\end{align}
The last line completes a square, retaining the curvature term which a
plain perturbation-norm estimate discards. Also,
\begin{equation}
 \tfrac12|\langle E,P_\star-P\rangle|
 \le v\sqrt{h_C(P)/\Delta_0}
 \le a h_C(P)+\frac{v^2}{4a\Delta_0}.
\end{equation}
For sharpness take $r=1,d=2$,
$C=\gamma I+\operatorname{Diag}(\Delta_0,0)$, $\gamma>0$,
$E_e=\left(\begin{smallmatrix}0&e\\e&0\end{smallmatrix}\right)$,
and $P_e=u_eu_e^\top$ with $u_e=(\cos\theta_e,\sin\theta_e)$,
$\theta_e=-e/(a\Delta_0)$. Direct diagonalization gives
\begin{equation}
 \lim_{e\to0}\frac{\Delta_0\{h_{C+E_e}(P_e)-(1+a)h_C(P_e)\}}
                         {\|E_e\|_F^2}=\frac{1+a^{-1}}4.
\end{equation}
\end{proof}

\begin{theorem}[Tracking with unrounded contraction and perturbation constants]
\label{thm:tracking}
Under Assumption~\ref{ass:tracking}, let
$\mu=\Delta_0\cos^2R_0$, $\alpha_t=\rho_t\mu$,
$L_R\rho_t\le1/4$, and $\alpha_t\le1/2$. Define
\begin{align}
 q_t&=(1+\alpha_t/2)(1-\alpha_t),\\
 A_t&=2(1+\alpha_t/2)
                 (L_R\rho_t^2\sigma_t^2+\rho_tB_t^2)
                  +\left(\frac1{2\alpha_t}+\frac14\right)
                                      \frac{v_t^2}{\Delta_0}.
\end{align}
For $h_t=h_{C_t}(U_t)$,
\begin{align}
 \mathbb E h_{t+1}&\le q_t\mathbb E h_t+A_t,
 \label{eq:tracking-recursion}\\
 \mathbb E h_T&\le h_0\prod_{j<T}q_j+
                  \sum_{t<T}A_t\prod_{j=t+1}^{T-1}q_j.
 \label{eq:tracking-unrolled}
\end{align}
For constant bounds and step $\rho$, with
$\alpha=\rho\mu$ and $\chi_\alpha=(1+\alpha/2)/(1+\alpha)$,
\begin{equation}
 \limsup_T\mathbb E h_T\le\chi_\alpha\left[
       \frac{4L_R\rho\sigma^2}{\mu}+\frac{4B^2}{\mu}
                      +\frac{v^2}{\rho^2\mu^2\Delta_0}\right].
 \label{eq:tracking-floor}
\end{equation}
Moreover $\mathbb E d_{\rm Gr}(U_T,U_T^\star)^2
\le\pi^2\mathbb Eh_T/(2\Delta_0)$.
\end{theorem}
\begin{proof}
For $h_t^+=h_{C_t}(U_{t+1})$, sharp QR pullback smoothness,
$\mathbb E_t\zeta_t=0$, and Young's inequality give
\begin{align}
 \mathbb E_t h_t^+
 &\le h_t-\rho_t\|g_t\|^2
       -\rho_t\langle g_t,\mathbb E_tb_t\rangle
       +\tfrac12L_R\rho_t^2\mathbb E_t\|g_t+\zeta_t+b_t\|^2\\
 &\le h_t-\tfrac12\rho_t\|g_t\|^2
                          +2L_R\rho_t^2\sigma_t^2+2\rho_tB_t^2\\
 &\le(1-\alpha_t)h_t+2L_R\rho_t^2\sigma_t^2+2\rho_tB_t^2.
 \label{eq:frozen-tracking}
\end{align}
Here $\rho_t|\langle g_t,\mathbb E_tb_t\rangle|
\le\rho_t\|g_t\|^2/4+\rho_tB_t^2$ and
$\mathbb E_t\|g_t+\zeta_t+b_t\|^2
\le2\|g_t\|^2+4\sigma_t^2+4B_t^2$; the step condition gives the
second line. Apply Lemma~\ref{lem:excess-drift} pathwise at $U_{t+1}$
with $a=\alpha_t/2$ and keep its constants to obtain the stated recursion.
Unroll it and use $1-q=\alpha(1+\alpha)/2$ for the steady bound.
The angle bound follows from quadratic growth and
$\sin\theta\ge2\theta/\pi$. Only the atomic smoothness and excess-drift
constants are claimed sharp; the stochastic bias budget is sufficient.
\end{proof}

The noise term increases with $\rho$, whereas the drift term decreases
with $\rho^2$. For the noise-and-drift terms, dropping
$\chi_\alpha\le1$ gives the envelope $A\rho+D/\rho^2$, whose minimizer is $\rho=(2D/A)^{1/3}$, subject to the
admissible-step restrictions. This is an explicit tracking trade-off, not a
claim that larger drift always increases GrassRoute's advantage over SVD.
If covariance changes are driven by the parameter update and
$\|C(W')-C(W)\|\le L_C\|W'-W\|$, then
\begin{equation}
 v_t\le L_C\eta_t\|D_t\|.
 \label{eq:coupling-drift}
\end{equation}
With bounded directions, a sufficient scaling for the drift contribution
in \eqref{eq:tracking-floor} to be $O(T^{-1/2})$ is
$\rho=O(T^{-1/2})$ and $v=O(T^{-3/4})$. The weaker assertion
$v=o(T^{-1/2})$ only implies a vanishing drift floor, without that rate.
Different schedules can trade parameter speed for tracking accuracy.

\paragraph{Momentum and retraction have separate costs.}
For the implemented update,
$H_t=\mu_H\Pi_{U_t}H_{t-1}+(1-\mu_H)\widehat Q_t$.
Writing $\widehat Q_t=Q_t+\epsilon_t$ and subtracting $Q_t$ gives
\begin{equation}
 H_t-Q_t=\mu_H\Pi_{U_t}(H_{t-1}-Q_{t-1})
 +\mu_H(\Pi_{U_t}Q_{t-1}-Q_t)+(1-\mu_H)\epsilon_t.
 \label{eq:momentum-lag}
\end{equation}
Minkowski's inequality and nonexpansiveness imply an $L_2$ lag bound
$e_t\le\mu_He_{t-1}+\mu_HD_t+(1-\mu_H)\sigma_t$, where
$D_t$ bounds the transported population-direction change. In a stationary
region with $D_t\le D$, this gives
$\limsup e_t\le\mu_HD/(1-\mu_H)+\sigma$.
The familiar variance reduction factor $(1-\mu_H)/(1+\mu_H)$ holds for
an EMA of independent noise at a \emph{frozen} basis; it cannot be assumed
for moving, data-dependent bases. Equation~\eqref{eq:momentum-lag}
states exactly which additional terms must be controlled.

If QR is compared with an exponential step of the \emph{same} tangent
direction, Lemma~\ref{lem:qr} gives a local $O(\rho_t^3\|H_t\|^3)$
subspace difference. When QR is analyzed directly through
\eqref{eq:pullback-smooth}, that error is already accounted for in the
pullback constant; adding it again as independent noise would double count.
Neither the curvature upper bound nor positive curvature alone supplies a
universal signed correction to the tracking recursion.

\paragraph{Continuous motion and jumps.}
Davis--Kahan also gives
$s(U_{t+1}^\star,U_t^\star)\le\sqrt2v_t/\Delta_0$.
A small covariance increment therefore bounds target motion. The converse
does not hold, since eigenvalues can change without rotating the target.
Theorem~\ref{thm:tracking} uses $v_t$ as a conservative observable control.
An abrupt change that violates the angle neighborhood invalidates the local
guarantee; it need not cause divergence, and no deterministic ordering
between online tracking and periodic recomputation follows. Boundary
experiments should test the effect of jump size while holding the spectrum
and gradient noise fixed.
